\documentclass[11pt,a4paper]{article}
\usepackage[utf8]{inputenc}
\usepackage[spanish]{babel}
\usepackage{geometry}
\geometry{top=2.5cm, bottom=2.5cm, left=3cm, right=3cm}
\usepackage{hyperref}
\usepackage{enumitem}
\usepackage{graphicx}
\usepackage{xcolor}
\usepackage{tcolorbox}

\definecolor{avisocolor}{RGB}{255, 243, 224}

% Captura: si la imagen aún no existe, se muestra un recuadro con lo que debe contener.
\newcommand{\captura}[2]{%
	\begin{center}
		\IfFileExists{img/#1}%
		{\includegraphics[width=0.75\textwidth]{img/#1}}%
		{\fbox{\parbox{0.85\textwidth}{\centering\vspace{1.5cm}\textit{Captura pendiente: #2}\vspace{1.5cm}}}}%
	\end{center}}

\title{Automatizar OSF con un asistente de IA \\ \large Guía rápida para no técnicos}
\author{Iván Rodríguez Steuerberg}
\date{Octubre de 2026}

\begin{document}
	\maketitle

	\section*{Qué vas a conseguir}
	Que un asistente de IA publique tus documentos en tu proyecto de OSF sin que tengas que subirlos a mano, conservando todas las versiones anteriores, y que mantenga un índice actualizado en la wiki del proyecto. Cómo organices tus carpetas es cosa tuya: esta guía solo cubre la conexión y la automatización.

	\section*{Qué necesitas}
	\begin{itemize}
		\item Una cuenta de OSF con tu proyecto creado.
		\item Un asistente de IA que pueda ejecutar órdenes en tu ordenador, como Claude Code o Antigravity. Un chat normal en el navegador no sirve, porque no puede acceder a tus archivos.
		\item Unos 15 minutos.
	\end{itemize}

	\section{Crear tu token de OSF}
	El token es una contraseña especial que permite al asistente trabajar con tu proyecto sin conocer tu contraseña de verdad.
	\begin{enumerate}[itemsep=4pt]
		\item En el menú lateral izquierdo de OSF, despliega \textbf{Settings} y entra en \textbf{Personal access tokens} (también puedes ir directamente a \url{https://osf.io/settings/tokens}).
		\item Pulsa \textbf{Create Token}.
	\end{enumerate}
	\captura{01-tokens.png}{página Personal access tokens de los ajustes de OSF}
	\begin{enumerate}[resume, itemsep=4pt]
		\item Ponle un nombre que reconozcas, por ejemplo ``Asistente IA''.
		\item Marca solo estos dos permisos: \texttt{osf.full\_read} (leer) y \texttt{osf.full\_write} (escribir).
		\item Pulsa \textbf{Create Token} y \textbf{copia el token en ese momento}: OSF no te lo volverá a enseñar.
	\end{enumerate}
	\captura{02-crear-token.png}{formulario de creación del token con los dos permisos marcados (sin ningún token visible)}

	\section{Guardar el token}
	\begin{enumerate}[itemsep=4pt]
		\item En la carpeta donde tienes tus documentos, crea un archivo de texto llamado exactamente \texttt{.env} (con el punto delante y sin \texttt{.txt} al final).
		\item Escribe dentro una sola línea: \texttt{OSF\_TOKEN=} seguido del token que copiaste.
		\item Guarda el archivo.
	\end{enumerate}

	\begin{tcolorbox}[colback=avisocolor, colframe=orange!70!black, title=Cuida tu token]
		\begin{itemize}
			\item No lo pegues nunca en un chat, ni siquiera en el del asistente: él lo leerá del archivo.
			\item No lo subas a GitHub ni lo compartas con nadie.
			\item Si crees que se ha filtrado, bórralo en la misma página de OSF y crea otro.
		\end{itemize}
	\end{tcolorbox}

	\section{Averiguar el identificador de tu proyecto}
	Abre tu proyecto en OSF y mira la barra de direcciones del navegador: el identificador son los cinco caracteres que van después de \texttt{osf.io/}. Por ejemplo, en la dirección \url{https://osf.io/ysmkz/overview} el identificador es \texttt{ysmkz}.

	\section{Pegar el prompt maestro}
	\begin{enumerate}[itemsep=4pt]
		\item Abre tu asistente de IA en la carpeta de tus documentos.
		\item Copia el prompt maestro del archivo \href{https://github.com/IvanSteuerberg/smart-basket/blob/main/docs/guia-osf/prompt-maestro.md}{\texttt{prompt-maestro.md}}, que está en GitHub, en el repositorio \textit{smart-basket} de IvanSteuerberg (carpeta \texttt{docs/guia-osf}).
		\item Cambia los datos entre \texttt{<...>} (el identificador del proyecto y la carpeta) y envíalo.
		\item El asistente primero solo leerá tu proyecto y te enseñará lo que encuentra. Después preparará los scripts y hará una simulación. Revisa lo que te propone y dile que sí solo si estás de acuerdo.
	\end{enumerate}

	\section{En el día a día}
	Basta con pedírselo:
	\begin{itemize}
		\item \textbf{``Publica mis documentos en OSF''}: te enseña la simulación y, si aceptas, sube solo lo que ha cambiado. OSF guarda la versión anterior de cada archivo, así que no hace falta poner fechas en los nombres.
		\item \textbf{``Actualiza el índice de la wiki de OSF''}: regenera la página principal con tus documentos, sus fechas y sus versiones.
	\end{itemize}

\end{document}
